% =============================================================================
\chapter{Results}
\label{ch:results}
% =============================================================================

% TODO: Structure
%
% \section{Text Model Results (H1)}
% - Coefficients for $\Delta$Sentiment, $\Delta$Risk, TextSim
% - Incremental $R^2$ over baseline
% - Expected: $\Delta$Risk positive, $\Delta$Sentiment negative, TextSim negative?
%
% \section{Divergence Model Results (H2)}
% - Coefficient for Divergence (the key result)
% - Expected: positive --- narrative-financial mismatch → higher future vol
% - Incremental $R^2$ over the text model
% - F-test for divergence significance
%
% \section{Economic Magnitude}
% - Portfolio sorts by divergence quartile → spread in future volatility
% - What does a 1-SD increase in divergence imply for annual vol?


% ─────────────────────────────────────────────────────────────────────────────
\section{Descriptive Statistics}
\label{sec:descriptive}
% ─────────────────────────────────────────────────────────────────────────────

Table~\ref{tab:summary} reports summary statistics for the regression
sample.  The dependent variable, post-filing annualised volatility, has a
mean of about 0.33 and a standard deviation of 0.17, showing large
cross-sectional differences in firm-level risk.  The interquartile range
(0.22 to 0.39) points to a right-skewed distribution, consistent with
occasional spikes during market stress periods such as 2020.

Lagged volatility has a nearly identical distribution, as expected for a
persistent series.  Firm size, measured as $\ln(\text{Total Assets})$,
ranges from mid-sized (P25~$\approx$~22.1) to large companies
(P75~$\approx$~24.2).  Leverage centres around 0.61 with moderate spread.
ROA has a mean near 0.06, and asset growth averages about 0.12 with a
long right tail driven by acquisitive firms.

\begin{table}[htbp]
\centering
\caption{Summary Statistics}
\label{tab:summary}
\begin{tabular}{lccccc}
\toprule
 & Mean & SD & P25 & Median & P75 \\
\midrule
$\ln(\sigma_{i,t+1})$ & -1.344 & 0.588 & -1.740 & -1.410 & -1.058 \\
$\ln(\sigma_{i,t})$ & -1.353 & 0.590 & -1.749 & -1.424 & -1.065 \\
$\ln(\text{Total Assets})$ & 23.182 & 1.452 & 22.191 & 23.206 & 24.236 \\
Leverage & 0.616 & 0.218 & 0.476 & 0.611 & 0.739 \\
ROA & 0.063 & 0.087 & 0.027 & 0.062 & 0.106 \\
Asset Growth & 0.113 & 0.243 & 0.001 & 0.059 & 0.142 \\
$\ln(\text{Words}^{\text{1A}})$ & 9.026 & 0.637 & 8.639 & 9.079 & 9.418 \\
RiskDensity & 0.051 & 0.006 & 0.047 & 0.051 & 0.055 \\
UncDensity & 0.034 & 0.005 & 0.032 & 0.034 & 0.037 \\
LitDensity & 0.017 & 0.005 & 0.013 & 0.016 & 0.019 \\
\midrule
Observations & \multicolumn{5}{c}{5,691} \\
Firms & \multicolumn{5}{c}{535} \\
\bottomrule
\end{tabular}
\par\smallskip
\noindent\footnotesize 
\textit{Note:} Summary statistics for the regression sample at the 30-day horizon. $\sigma_{i,t+1}$ is the annualised standard deviation of daily log returns over the 30-day window starting two trading days after the filing date; $\ln(\sigma_{i,t+1})$ is its natural logarithm. Words$^{\text{1A}}$ is the cleaned word count of Item~1A; RiskDensity is the share of LM uncertainty- and litigious-list words in Item~1A; UncDensity and LitDensity are the two sub-components separately. Financial ratios are winsorised at the 1st and 99th percentiles.
\end{table}

Table~\ref{tab:correlation} presents pairwise Pearson correlations.
Volatility and its lag are strongly positively correlated ($r = 0.61$),
confirming persistence.  Firm size correlates negatively with volatility
($r \approx -0.32$), in line with the well-known size--risk relationship.
ROA also correlates negatively with volatility ($r \approx -0.33$),
while leverage and asset growth show weaker univariate links.
None of the pairwise correlations among the explanatory variables
exceed 0.40 in absolute value, reducing multicollinearity concerns
(variance inflation factors are below 1.4 for all regressors).

\begin{table}[htbp]
\centering
\caption{Pearson Correlation Matrix}
\label{tab:correlation}
\begin{tabular}{lcccccccccc}
\toprule
 & (1) & (2) & (3) & (4) & (5) & (6) & (7) & (8) & (9) & (10) \\
\midrule
(1) $\ln(\sigma_{i,t+1})$ & 1 &  &  &  &  &  &  &  &  &  \\
(2) $\ln(\sigma_{i,t})$ & 0.35*** & 1 &  &  &  &  &  &  &  &  \\
(3) $\ln(\text{Total Assets})$ & -0.19*** & -0.20*** & 1 &  &  &  &  &  &  &  \\
(4) Leverage & -0.06*** & -0.07*** & 0.21*** & 1 &  &  &  &  &  &  \\
(5) ROA & -0.23*** & -0.25*** & 0.05*** & -0.06*** & 1 &  &  &  &  &  \\
(6) Asset Growth & 0.16*** & 0.09*** & -0.12*** & -0.12*** & -0.04*** & 1 &  &  &  &  \\
(7) $\ln(\text{Words}^{\text{1A}})$ & 0.31*** & 0.31*** & -0.09*** & -0.04*** & -0.27*** & 0.17*** & 1 &  &  &  \\
(8) RiskDensity & -0.11*** & -0.12*** & -0.01 & -0.01 & 0.06*** & -0.08*** & -0.30*** & 1 &  &  \\
(9) UncDensity & -0.09*** & -0.08*** & -0.18*** & -0.06*** & 0.14*** & -0.05*** & -0.43*** & 0.71*** & 1 &  \\
(10) LitDensity & -0.06*** & -0.08*** & 0.18*** & 0.04*** & -0.06*** & -0.06*** & 0.03** & 0.66*** & -0.06*** & 1 \\
\bottomrule
\end{tabular}
\par\smallskip
\noindent\footnotesize 
\textit{Note:} Pairwise Pearson correlation coefficients on the 30-day-horizon regression sample. ***, **, * denote significance at the 1\%, 5\%, and 10\% levels.
\end{table}


% ─────────────────────────────────────────────────────────────────────────────
\section{Baseline: Financial Determinants of Volatility}
\label{sec:baseline}
% ─────────────────────────────────────────────────────────────────────────────

Table~\ref{tab:baseline} reports OLS estimates of post-filing
annualised volatility on progressively richer sets of financial
controls, estimated as described in Section~\ref{sec:estimation}.
The three columns correspond to nested specifications: lagged
volatility only (Column~1), with firm size and leverage added
(Column~2), and the full baseline including return on assets and
asset growth (Column~3).

\begin{table}[htbp]
\centering
\caption{Baseline Volatility Determinants}
\label{tab:baseline}
\begin{tabular}{lccc}
\toprule
 & (1) & (2) & (3) \\
\midrule
Lagged Volatility & $+0.6697$*** & $+0.6108$*** & $+0.5483$*** \\
 & $(31.98)$ & $(29.54)$ & $(19.97)$ \\[4pt]
$\ln(\text{Total Assets})$ &  & $-0.0163$*** & $-0.0162$*** \\
 &  & $(-7.26)$ & $(-7.36)$ \\[4pt]
Leverage &  & $+0.0029$ & $+0.0004$ \\
 &  & $(0.29)$ & $(0.04)$ \\[4pt]
ROA &  &  & $-0.2260$*** \\
 &  &  & $(-8.01)$ \\[4pt]
Asset Growth &  &  & $+0.0267$*** \\
 &  &  & $(3.24)$ \\[4pt]
\midrule
Industry FE & Yes & Yes & Yes \\
Year FE & Yes & Yes & Yes \\
Observations & 6,349 & 6,349 & 6,349 \\
Firms & 540 & 540 & 540 \\
Adj.\ $R^2$ & 0.626 & 0.639 & 0.651 \\
\bottomrule
\end{tabular}
\begin{tablenotes}
\item \textit{Note:} This table reports OLS estimates of post-filing annualised volatility on financial determinants. The dependent variable is the annualised standard deviation of daily log returns over the 365-day window following each 10-K filing date. All specifications include two-digit SIC industry (49 groups) and fiscal-year fixed effects. $t$-statistics, reported in parentheses, are based on standard errors clustered at the firm level. ***, **, and * denote significance at the 1\%, 5\%, and 10\% levels, respectively.
\end{tablenotes}
\end{table}

\paragraph{Volatility persistence.}
Column~(1) includes only lagged volatility ($\sigma_{i,t}$) alongside the
fixed effects.  The coefficient of $0.67$ ($t = 33.5$) confirms the
well-documented persistence of firm-level return volatility: a one
percentage-point rise in last year's annualised volatility maps into
about a two-thirds percentage-point rise the following year.  This single
regressor, combined with industry and year fixed effects, already
explains about $63\%$ of the cross-sectional and time-series variation
in post-filing volatility, showing that most of the predictable component
of risk is mechanical.

\paragraph{Firm size and capital structure.}
Column~(2) adds firm size, measured as the natural logarithm of total
assets, and book leverage (total liabilities over total assets).  Size
enters with a negative and highly significant coefficient
($-0.017$, $t = -7.4$): larger firms exhibit lower future volatility,
consistent with the size--risk relationship documented for U.S.\ equities
by \textcite{chi2017} and \textcite{campbell2014}.  Leverage is
economically negligible and statistically indistinguishable from zero
($t = 0.17$); the industry fixed effects appear to absorb most of the
cross-sectional dispersion in capital structure, leaving little residual
variation through which leverage can operate.

\paragraph{Profitability and growth.}
The full specification in Column~(3) adds return on assets and asset
growth.  ROA carries a large negative coefficient ($-0.224$,
$t = -7.9$), confirming that profitable firms face lower subsequent
volatility, in line with the earnings-quality intuition of
\textcite{francis2005}.  Asset growth enters positively
($+0.026$, $t = 3.1$): rapid expansion is associated with elevated
future uncertainty, consistent with the operational-risk channel.
Adding these two regressors raises the adjusted~$R^2$ from $0.643$ to
$0.654$---a modest but statistically meaningful improvement that refines,
rather than dominates, the persistence-and-size benchmark.

A noteworthy by-product of the full specification is that the lagged
volatility coefficient declines from $0.67$ in Column~(1) to $0.55$
($t = 20.7$) in Column~(3).  Roughly one-fifth of what looked like raw
persistence is in fact mediated by slower-moving fundamentals---size,
profitability, and growth---that are themselves persistent.  Lagged
volatility nonetheless remains by a wide margin the most informative
single regressor in the model.

\paragraph{Benchmark for textual analysis.}
The full baseline model explains $65.4\%$ of the variation in
post-filing volatility.  This adjusted~$R^2$ is the bar that the
text model and the divergence model must clear, on the criteria
listed in Section~\ref{sec:hypothesis_tests}.  The bar is high:
financial fundamentals together with volatility persistence already
capture most of the predictable variation in firm-level risk,
leaving narrative disclosures to compete for what remains.
